\documentclass[10pt,a4paper]{amsart}
\usepackage{amsmath,amssymb,amsthm,graphicx}
% For commutative diagrams we use Xy-pic
\usepackage[all]{xy}

% We use 2cell for 2-commutative diagrams.
\xyoption{2cell}
% Theorem environments.
%
\theoremstyle{plain}
\newtheorem{theorem}[subsection]{Theorem}
\newtheorem{proposition}[subsection]{Proposition}
\newtheorem{lemma}[subsection]{Lemma}

\theoremstyle{definition}
\newtheorem{definition}[subsection]{Definition}
\newtheorem{example}[subsection]{Example}
\newtheorem{exercise}[subsection]{Exercise}
\newtheorem{situation}[subsection]{Situation}

\theoremstyle{remark}
\newtheorem{remark}[subsection]{Remark}
\newtheorem{remarks}[subsection]{Remarks}

\numberwithin{equation}{subsection}

% Macros
%
\def\lim{\mathop{\mathrm{lim}}\nolimits}
\def\colim{\mathop{\mathrm{colim}}\nolimits}
\def\Spec{\mathop{\mathrm{Spec}}}
\def\Hom{\mathop{\mathrm{Hom}}\nolimits}
\def\Ext{\mathop{\mathrm{Ext}}\nolimits}
\def\SheafHom{\mathop{\mathcal{H}\!\mathit{om}}\nolimits}
\def\SheafExt{\mathop{\mathcal{E}\!\mathit{xt}}\nolimits}
\def\Sch{\mathit{Sch}}
\def\Mor{\mathop{\mathrm{Mor}}\nolimits}
\def\Ob{\mathop{\mathrm{Ob}}\nolimits}
\def\Sh{\mathop{\mathit{Sh}}\nolimits}
\def\NL{\mathop{N\!L}\nolimits}
\def\CH{\mathop{\mathrm{CH}}\nolimits}
\def\proetale{{pro\text{-}\acute{e}tale}}
\def\etale{{\acute{e}tale}}
\def\QCoh{\mathit{QCoh}}
\def\Ker{\mathop{\mathrm{Ker}}}
\def\Im{\mathop{\mathrm{Im}}}
\def\Coker{\mathop{\mathrm{Coker}}}
\def\Coim{\mathop{\mathrm{Coim}}}

% Boxtimes
%
\DeclareMathSymbol{\boxtimes}{\mathbin}{AMSa}{"02}

%
% Macros for moduli stacks/spaces
%
\def\QCohstack{\mathcal{QC}\!\mathit{oh}}
\def\Cohstack{\mathcal{C}\!\mathit{oh}}
\def\Spacesstack{\mathcal{S}\!\mathit{paces}}
\def\Quotfunctor{\mathrm{Quot}}
\def\Hilbfunctor{\mathrm{Hilb}}
\def\Curvesstack{\mathcal{C}\!\mathit{urves}}
\def\Polarizedstack{\mathcal{P}\!\mathit{olarized}}
\def\Complexesstack{\mathcal{C}\!\mathit{omplexes}}
% \Pic is the operator that assigns to X its picard group, usage \Pic(X)
% \Picardstack_{X/B} denotes the Picard stack of X over B
% \Picardfunctor_{X/B} denotes the Picard functor of X over B
\def\Pic{\mathop{\mathrm{Pic}}\nolimits}
\def\Picardstack{\mathcal{P}\!\mathit{ic}}
\def\Picardfunctor{\mathrm{Pic}}
\def\Deformationcategory{\mathcal{D}\!\mathit{ef}}
\usepackage{cleveref}
\crefname{originallemma}{original lemma}{original lemmas}
\Crefname{originallemma}{Original lemma}{Original lemmas}
\makeatletter
\newlabel{lemma-closed-immersion-almost-essentially-surjective}{{04FW}{source}{}{sourcepointer-1}{}}
\newlabel{lemma-closed-immersion-almost-essentially-surjective@cref}{{[originallemma][0][]04FW}{[1][0][]source}{}{}{}}
\makeatother
\crefname{originallemma}{original lemma}{original lemmas}
\Crefname{originallemma}{Original lemma}{Original lemmas}
\makeatletter
\newlabel{lemma-closed-immersion-almost-full}{{04FV}{source}{}{sourcepointer-2}{}}
\newlabel{lemma-closed-immersion-almost-full@cref}{{[originallemma][0][]04FV}{[1][0][]source}{}{}{}}
\makeatother
\crefname{originallemma}{original lemma}{original lemmas}
\Crefname{originallemma}{Original lemma}{Original lemmas}
\makeatletter
\newlabel{categories-definition-cofinal}{{04E6}{source}{}{sourcepointer-3}{}}
\newlabel{categories-definition-cofinal@cref}{{[originallemma][0][]04E6}{[1][0][]source}{}{}{}}
\makeatother
\crefname{originallemma}{original lemma}{original lemmas}
\Crefname{originallemma}{Original lemma}{Original lemmas}
\makeatletter
\newlabel{categories-lemma-cofinal}{{04E7}{source}{}{sourcepointer-4}{}}
\newlabel{categories-lemma-cofinal@cref}{{[originallemma][0][]04E7}{[1][0][]source}{}{}{}}
\makeatother
\title{Stalks of the pushforward along a closed immersion}
\author{The Stacks Project Authors}
\date{stacks-project:etale-cohomology.tex:local-snapshot-20261011; GNU FDL 1.2 or later}
\begin{document}
\maketitle
\noindent\textit{Editorial note.} All delivered mathematical text is the author's own, copied from the pinned Stacks Project chapter file etale-cohomology.tex; no statement or proof step was rewritten, reordered or removed.\par\medskip
\noindent\textit{Editorial note.} The proof is the author's complete argument for both cases of the displayed formula.\par\medskip
\noindent\textit{Editorial note.} Every labelled reference occurring inside the counted spans is carried in assembly.external\_references as a non-load-bearing navigation pointer (same-chapter ones at their definition site, cross-chapter ones at the invocation inside the counted span), so no load-bearing mathematics is replaced by a stub. No bibliography citation occurs inside the counted spans.\par\medskip
\noindent\textit{Editorial note.} The checked chapter list for same-name collisions is the set of stacks chapters staged in this repository: chow, cohomology, divisors, etale-cohomology, flat, more-algebra, more-morphisms, morphisms, spaces-more-morphisms, varieties together with this card's chapter etale-cohomology.\par\medskip
\noindent\textit{Editorial note.} The preamble is the chapter's own native preamble excerpt already used for this source: the Xy-pic setup (native preamble lines 14-18) and the native theorem declarations with the mathematical macros (native preamble lines 154-217).\par\medskip
\noindent\textit{Allowed standard prerequisites (editorial).}\begin{itemize}
\item Closed immersion i : Z -> X and the open complement U = X minus Z. Same chapter, the open immersion of the complement and the fact that the small etale site of X is covered by Z and U. Supplies the two cases of the proof: a geometric point inside Z and one outside.
\item Stalk of a sheaf at a geometric point and the description of the pushforward by sections. Same chapter, the stalk functor at a geometric point and the description of sections of i\_\{small,*\}G. Supplies the computation in both cases.
\end{itemize}\medskip
\ifdefined\hypertarget\hypertarget{sourcepointer-1}{}\fi
\noindent\textit{Source reference.} Original lemma 04FW: sources/stacks-etale-cohomology/pinned-original-source.tex, source lines 6621--6621. Navigation pointer invoked at proof line 6693; the counted span contains the pointer only and the cited result is not reproduced inside it, so the pointer is registered as a non-load-bearing navigation stub in the established form for Stacks material. No bibliography citation is involved.\par\medskip
\ifdefined\hypertarget\hypertarget{sourcepointer-2}{}\fi
\noindent\textit{Source reference.} Original lemma 04FV: sources/stacks-etale-cohomology/pinned-original-source.tex, source lines 6597--6597. Navigation pointer invoked at proof line 6694; the counted span contains the pointer only and the cited result is not reproduced inside it, so the pointer is registered as a non-load-bearing navigation stub in the established form for Stacks material. No bibliography citation is involved.\par\medskip
\ifdefined\hypertarget\hypertarget{sourcepointer-3}{}\fi
\noindent\textit{Source reference.} Original lemma 04E6: sources/stacks-etale-cohomology/pinned-original-source.tex, source lines 6696--6696. Navigation pointer invoked at proof line 6696; the counted span contains the pointer only and the cited result is not reproduced inside it, so the pointer is registered as a non-load-bearing navigation stub in the established form for Stacks material. No bibliography citation is involved.\par\medskip
\ifdefined\hypertarget\hypertarget{sourcepointer-4}{}\fi
\noindent\textit{Source reference.} Original lemma 04E7: sources/stacks-etale-cohomology/pinned-original-source.tex, source lines 6698--6698. Navigation pointer invoked at proof line 6698; the counted span contains the pointer only and the cited result is not reproduced inside it, so the pointer is registered as a non-load-bearing navigation stub in the established form for Stacks material. No bibliography citation is involved.\par\medskip
\begin{lemma}
\label{lemma-stalk-pushforward-closed-immersion}
Let $i : Z \to X$ be a closed immersion of schemes.
Let $\mathcal{G}$ be a sheaf of sets on $Z_\etale$.
Let $\overline{x}$ be a geometric point of $X$.
Then
$$
(i_{small, *}\mathcal{G})_{\overline{x}} =
\left\{
\begin{matrix}
* & \text{if} & \overline{x} \not \in Z \\
\mathcal{G}_{\overline{x}} & \text{if} & \overline{x} \in Z
\end{matrix}
\right.
$$
where $*$ denotes a singleton set.
\end{lemma}

\begin{proof}
Let $U = X \setminus Z$.
Note that $i_{small, *}\mathcal{G}|_{U_\etale} = *$ is the final
object in the category of \'etale sheaves on $U$, i.e., the sheaf
which associates a singleton set to each scheme \'etale over $U$.
This explains the value of $(i_{small, *}\mathcal{G})_{\overline{x}}$
if $\overline{x} \not \in Z$.

\medskip\noindent
Next, suppose that $\overline{x} \in Z$. Note that
$$
(i_{small, *}\mathcal{G})_{\overline{x}}
=
\colim_{(U, \overline{u})} \mathcal{G}(U_Z)
$$
and on the other hand
$$
\mathcal{G}_{\overline{x}}
=
\colim_{(V, \overline{v})} \mathcal{G}(V).
$$
Let $\mathcal{C}_1 = \{(U, \overline{u})\}^{opp}$ be the opposite of the
category of \'etale neighbourhoods of $\overline{x}$ in $X$, and let
$\mathcal{C}_2 = \{(V, \overline{v})\}^{opp}$ be the opposite of the
category of \'etale neighbourhoods of $\overline{x}$ in $Z$. The canonical map
$$
\mathcal{G}_{\overline{x}}
\longrightarrow
(i_{small, *}\mathcal{G})_{\overline{x}}
$$
corresponds to the functor $F : \mathcal{C}_1 \to \mathcal{C}_2$,
$F(U, \overline{u}) = (U_Z, \overline{x})$. Now
Lemmas \ref{lemma-closed-immersion-almost-essentially-surjective} and
\ref{lemma-closed-immersion-almost-full}
imply that $\mathcal{C}_1$ is cofinal in $\mathcal{C}_2$, see
Categories, Definition \ref{categories-definition-cofinal}.
Hence it follows that the displayed arrow is an isomorphism, see
Categories, Lemma \ref{categories-lemma-cofinal}.
\end{proof}

\end{document}
